\section{Background geometry and finite branch tubes}\label{sec:background}

This section states the precise background inputs and derives the
geometric properties of the finite observation tubes. The background
neighborhood and branch-tube choices are made before any higher
derivative order is selected.

\subsection{Normalization of a fixed Kerr center}\label{bg:normalization}

The companion interfaces use mass one and positive rotation. For
$\mathfrak a<0$, the fixed axial reflection $\varphi\mapsto-\varphi$,
with Killing time unchanged, identifies the Kerr metric with parameter
$\mathfrak a$ with that for $|\mathfrak a|$; see the Boyer--Lindquist
metric in \cite[Section~7.1]{KerrCTwo}. It carries the bridge and its
future normal to the corresponding bridge. Simultaneous pullback of
the data and reference geometry preserves every covariant derivative
order, as in \cite[Section~2]{KerrCOne}.

After the reflection, reset notation for the reflected original bridge.
Let $\Psi:\Sigma_n\to\Sigma$ be the fixed identification from the
dimensionless normalized bridge coordinates to this bridge.
Constant spacetime scaling by $M^{-2}$ sends the data
increments to
\[
 \mathcal T(u,v)=(M^{-2}\Psi^*u,M^{-1}\Psi^*v)
\]
and the center to
$d_{*,n}=(M^{-2}\Psi^*h_*,M^{-1}\Psi^*K_*)$, the Kerr bridge with
parameters $(1,|\mathfrak a|/M)$. Write $h_{*,n}=M^{-2}\Psi^*h_*$.
Its connection is
the pullback of $\nabla_*$, since constant metric scaling does not
change the Levi-Civita connection. A covariant tensor of rank $2+j$
has norm scaled by $M^{2+j}$, so
\[
\begin{split}
 |\nabla_n^j(M^{-2}\Psi^*u)|_{h_{*,n}}
   &=M^j(|\nabla_*^ju|_*)\circ\Psi,\\
 |\nabla_n^j(M^{-1}\Psi^*v)|_{h_{*,n}}
   &=M^{j+1}(|\nabla_*^jv|_*)\circ\Psi .
\end{split}
\]
Let $p_m^{\rm orig}$ be \eqref{intro:seminorms} for $(h_*,\rho)$ and
$p_m^{\rm norm}$ its counterpart for $(h_{*,n},\rho_n)$, where $\rho_n$
is a fixed admissible normalized radial weight. The two weights are
comparable under $\Psi$. Consequently,
for each finite $m$ there is a fixed $C_m\ge1$ such that
\begin{equation}\label{bg:seminorm-scaling}
 C_m^{-1}p_m^{\rm orig}(u,v)
 \le p_m^{\rm norm}(\mathcal T(u,v))
 \le C_m p_m^{\rm orig}(u,v).
\end{equation}
This is the fixed-order comparison in the scaling discussions of
\cite{KerrCTwo,KerrCOne}; it shows in particular that the inverse image
of a normalized $p_{10}$ ball contains an original $p_{10}$ ball.

The fixed transformations preserve completeness, the vacuum equations,
maximality, and future orientation. They also preserve
Definition~\ref{intro:extension}: smooth coordinate changes preserve
the local Sobolev condition, and constant metric scaling leaves the
connection and timelike exits unchanged. The induced maps of weighted
constraint spaces are homeomorphisms by \eqref{bg:seminorm-scaling}.
Meagreness pulls back under these maps and restricts to the original
open ball. We may therefore prove the theorem below for $M=1$ and one
fixed $0<\mathfrak a<1$.

\subsection{The comparison development}

For the normalized fixed center, take a sufficiently
small compact neighborhood of physical Kerr parameters contained in the
rotating subextremal range. The exact-profile construction in
\cite[Section~7.1]{KerrCTwo} has $r_->0$ and positive absolute inner surface
gravity there. It rescales both optical directions to one fixed rate
$\kappa>0$, for example the center's absolute inner surface gravity, and
gives smooth positive relative lapse and sphere metric with smooth inverses
at the inner boundary. Choose its depth $q_0>0$ sufficiently small for both
angular gauges and for its exact-profile inequality
$(|\partial_\vartheta F|^2+\mathfrak a^2\sin^2\vartheta)_{\mathrm{entry}}
<4\mathfrak a^2$, where $F$ is the optical correction in that construction.
This is possible because the fixed
compact parameter range stays away from zero rotation. The deeper computing
exit, exterior constants, and data threshold are chosen afterward. All these
choices may depend on the fixed center; none is required to remain uniform
toward zero rotation or extremality.

Choose $\eps$ sufficiently small for the exterior and interior
constructions in \cite[\CompanionRef{foundation}{far:closure}, \CompanionRef{foundation}{in:entry}, \CompanionRef{foundation}{in:low}, \CompanionRef{foundation}{in:pure-jets}, \CompanionRef{foundation}{in:higher}, \CompanionRef{foundation}{in:continuation}, \CompanionRef{foundation}{cmp:development}]{KerrCTwo}.
After the finite bridge and profile preparation is fixed,
\CompanionRef{foundation}{far:closure} derives exterior closure from
\eqref{intro:seminorms}--\eqref{intro:neighborhood}, and
\CompanionRef{foundation}{in:entry} supplies the cylinder hypotheses
of the interior propositions. Their finite-slab and geometric
conditions are supplied by the same construction. The resulting
package requires no higher-seminorm smallness and no prescribed leading
coefficient or fixed-charge condition beyond the stated weighted bounds
and $p_{10}$ smallness. It gives a smooth vacuum comparison development
$\mathcal P_d$ with the entire bridge as Cauchy surface and an
open isometric embedding into $M_d$ fixing that surface. We
identify $\mathcal P_d$ with this image.

For clarity about dependencies, the companion uses the Kerr stationary
linear estimates and zero-mode classification in
\cite[Theorems~5.1, 5.2, 6.6, Proposition~6.8, and Lemma~6.9]{HHV2025}.
The trapping and system estimates used in
\cite[\CompanionRef{foundation}{sec:stationary}, \CompanionRef{foundation}{sec:compact}]{KerrCTwo} come from
\cite{Dyatlov2016,Hintz2017}. In that application, Hintz's
Condition~(2.5) controls the subprincipal margin and
Proposition~3.12 supplies the pseudodifferential-inner-product
reduction; the trapping hypotheses and the specialization to the
asymptotically flat Kerr problem are supplied in the companion.
Its finite-time nonlinear closure is a further argument there.
We use the resulting background statements through the interfaces below.

The comparison development contains a lower region with a past
collar, exterior computing regions on both ends, and a future
double-null wedge attached at a full spacelike cylinder
$\mathcal S\simeq\R\times\Sph$. In the gluing construction,
$\mathcal S$ is the wedge's only relative boundary inside
$\mathcal P_d$; see the proof of
\cite[\CompanionRef{foundation}{cmp:development}]{KerrCTwo}.
A future timelike track that crosses $\mathcal S$ therefore remains
in the wedge until its first runout from $\mathcal P_d$. In the wedge,
\begin{equation}\label{bg:metric}
\begin{split}
 t&=(u+v)/2\ge0,\qquad x=(v-u)/2,\qquad \mathcal S=\{t=0\},\\
 g&=-2a\,\dd u\,\dd v+
 \gamma_{AB}(\dd\theta^A-b^A\dd u)(\dd\theta^B-b^B\dd u),\\
 a&=q\mathcal A,\qquad q=q_0e^{-\kappa(u+v)},\qquad q_0,\kappa>0.
\end{split}
\end{equation}
The angular fields are tensors on $\Sph$. A fixed finite sphere
atlas is used for coordinate norms. Either branch can be arranged
to have $v\to+\infty$ by interchanging the eikonals from the
initial cylinder and preparing the angular labels for the new
shift-free direction. This is the label change described in
\cite[\CompanionRef{signed}{bg:geometry}]{KerrCOne} and
\cite[\CompanionRef{foundation}{in:label-swap}]{KerrCTwo}; it preserves the
following estimates.

Set $y_0=u,y_1=v$, $B_0=b,B_1=0$, and
\[
 L_i=\partial_{y_i}+B_i,
 \qquad \mathfrak L_i=\partial_{y_i}+\Lie_{B_i}.
\]
The second operator acts on angular tensors. The geometric fields
are
\begin{equation}\label{bg:fields}
\begin{gathered}
 \chi_i=\tfrac12\mathfrak L_i\gamma,
 \qquad \ell_i=L_i\log a,
 \qquad \tau_i=\tr_\gamma\chi_i,\\
 \xi_0+\xi_1=\dd_\theta\log a,
 \qquad \xi_0-\xi_1=-a^{-1}\gamma(\partial_vb,\cdot).
\end{gathered}
\end{equation}
Thus $\dd_\theta\log a=\dd_\theta\log\mathcal A$. Angular
covariant derivatives and contractions in the equations use $\gamma$
unless specified otherwise. Angular Sobolev $H^k$ norms use component
derivatives in the fixed finite atlas, with fixed localizers and
coordinate densities. Equivalently, one may use derivatives from one
fixed smooth reference connection and a fixed reference metric, with
solution-independent equivalence constants. This is the fixed-chart
convention of \cite[\CompanionRef{foundation}{in:low},
\CompanionRef{foundation}{in:higher}]{KerrCTwo}. In particular, the
metric and inverse bounds below give uniform pointwise comparison with
the reference metric.

\begin{convention}[Fixed-order subexponential bounds]
\label{conv:bg-subexp}
For a fixed datum and fixed finite derivative orders, a bound
$Q_T\le e^{o(T)}$ means that for every $\delta>0$ there are
$C_\delta,T_\delta$ such that
$Q_T\le C_\delta e^{\delta T}$ for $T\ge T_\delta$.
For a positive quantity, a two-sided bound also bounds its
reciprocal. Constants can depend on the datum and the chosen
orders. Uniformity for a family of changing data is asserted only
when proved separately.
\end{convention}

\begin{proposition}[Background estimates]\label{prop:bg-package}
For every $d\in\Ueps$, the comparison geometry has the following
properties.
\begin{enumerate}[label=\textup{(\roman*)}]
\item The metrics $\gamma^{\pm1}$ and $\mathcal A^{\pm1}$ are
bounded in angular $H^6$, and $\xi_i,a^{-1}\partial_vb$ are
bounded in angular $H^5$. There are nonnegative measurable functions
$g_i:\R\to[0,\infty)$ with
\begin{equation}\label{bg:cell-sum}
 \sum_{n\in\Z}\sup_{[n,n+1]}g_i<\infty
\end{equation}
such that
\begin{equation}\label{bg:directional-low}
 \norm{\chi_i}_{H^4}+\norm{\ell_i+\kappa}_{H^4}
       +\norm{B_i}_{H^5}
 \le C\bigl(q+g_i(y_i)\bigr).
\end{equation}
The corresponding first pure longitudinal derivatives
$\mathfrak L_i\chi_i,L_i\ell_i$, and $\partial_u b$ in direction
$0$, have directional summable bounds through angular $H^2$,
after enlarging the $g_i$.
\item The first ordinary longitudinal jets satisfy
\begin{equation}\label{bg:low-jets}
\begin{gathered}
 \norm b_{H^5}+\sum_i\norm{(\chi_i,\ell_i+\kappa)}_{H^4}\to0,\\
 \sum_{k=0}^1\left(\norm{\partial_{y_k}b}_{H^2}
       +\sum_i\norm{\partial_{y_k}(\chi_i,\ell_i)}_{H^2}\right)\to0,\\
 \sum_{k=0}^1\left(\norm{\partial_{y_k}\gamma}_{H^4}
       +\norm{\partial_{y_k}\log a+\kappa}_{H^4}\right)\to0
\end{gathered}
\end{equation}
as $t\to\infty$ and $|u|,|v|\to\infty$. The displayed norms
are bounded everywhere in the wedge.
\item For fixed $U\in\R$, $C>0$, and each fixed ordinary
coordinate derivative order, the fields
\begin{equation}\label{bg:high-fields}
 \gamma^{\pm1},\quad\log\mathcal A,\quad b,
       \quad a^{-1}\partial_vb
\end{equation}
have $e^{o(T)}$ bounds on
\begin{equation}\label{bg:truncated-wedge}
 u\le U,\qquad v\le CT,\qquad t\ge0.
\end{equation}
The assertion holds for either optical ordering and on restricted
coordinate-path segments. Fixed derivatives of $a$, $\sqrt a$,
and $\partial_vb$ retain respectively the factors $a$, $\sqrt a$,
and $a$, with subexponential coefficient bounds.
\end{enumerate}
\end{proposition}

\begin{proof}
These are \cite[\CompanionRef{foundation}{in:low}, \CompanionRef{foundation}{in:pure-jets}, \CompanionRef{foundation}{in:higher}]{KerrCTwo}, with cylinder
data supplied by \CompanionRef{foundation}{in:entry} and exterior closure supplied
by \CompanionRef{foundation}{far:closure} there; see also
\cite[\CompanionRef{signed}{bg:geometry}]{KerrCOne}. The low characteristic estimates
give \eqref{bg:directional-low} and summable cell suprema.
\CompanionRef{foundation}{in:pure-jets} includes the same-direction derivatives, which
are needed in \eqref{bg:low-jets}. After enlarging the envelopes to
include these jets, replace each envelope on $[n,n+1)$ by its
supremum on $[n,n+1]$. This gives a measurable step function whose
sum of closed-cell suprema increases by at most a factor two.
The frozen profile terms have
the stated $q$ bounds and are absorbed in the right sides.
\CompanionRef{foundation}{in:higher} provides \eqref{bg:high-fields} at every fixed
order. The last statement follows by differentiating
$a=q\mathcal A$ and the relative shift identity in
\eqref{bg:fields}. None of these results assumes a bound for
parallel-frame curvature on an observing geodesic.
\end{proof}

For the linear packet, one also uses the normalized vectors
\[
 \widehat e_i=a^{-1/2}L_i,\qquad \widehat e_A=\partial_A,
 \qquad Z_i=L_i,\qquad Z_A=\sqrt a\,\partial_A.
\]
Write $V_\alpha$ for the connection matrix of $\nabla_{Z_\alpha}$
in this frame. These matrices are bounded;
in the bulk $\min(t,v,-u)\to\infty$ their limits are
\begin{equation}\label{bg:normalized-limits}
 V_0\to D_0,\qquad V_1\to-D_0,\qquad V_A\to0,
 \qquad D_0=\operatorname{diag}(-\kappa/2,\kappa/2,0,0).
\end{equation}
The components of $a\Riem$ in normalized slots and the
differentiated connection arrays entering the wave operator are
bounded and tend to zero there; their fixed ordinary jets cost
$e^{o(T)}$. This is
\cite[\CompanionRef{signed}{bg:coefficients}]{KerrCOne}, with its explicit slot formulas,
equivalently \cite[\CompanionRef{foundation}{pk2:coefficients}]{KerrCTwo}. It is a consequence of
\Cref{prop:bg-package} and the vacuum identities, with no
bounded-curvature hypothesis on the track. This intermediate
input is the one used in the later linear transfer.

\subsection{Seeds, section measure, and first runout}

A seed $w\in T\Sigma$ determines a future unit initial velocity
whose normal and tangential components are
\[
 \bigl(\sqrt{1+h(w,w)},w\bigr).
\]
The maximal geodesic with that velocity is always understood in
the full $M_d$. Identify $T\Sigma$ with $T^*\Sigma$ by
$w\mapsto h(w,\cdot)$, and let $\mu_d$ be the pullback of its
canonical symplectic volume. In position and spatial covariant
momentum coordinates this is ordinary coordinate volume. It is
a smooth positive measure; on a fixed precompact seed chart it
depends continuously on the initial metric.

We will repeatedly use its section version. On the future unit
mass shell of the geodesic Hamiltonian, the restriction of the
symplectic form has the flow direction as kernel. A
single-valued transverse hitting map therefore preserves the
restricted two-form and its third exterior power. Thus the
measure on a time-level section is canonical position--momentum
volume, even when hitting time varies with the initial state.
This statement concerns smooth flow strictly inside $M_d$.

\begin{proposition}[First runout and corner removal]
\label{prop:bg-runout}
Every seed whose future geodesic has finite maximal proper
lifetime in $M_d$ crosses $\mathcal S$. Its first runout from
$\mathcal P_d$ is either a branch, where one optical coordinate
has a finite limit and the other tends to $+\infty$, or a corner,
where both tend to $+\infty$. For each fixed datum, seeds whose
first runout is a corner form a $\mu_d$-null set.
\end{proposition}

\begin{proof}
The full-lifetime classification is
\cite[\CompanionRef{foundation}{cmp:classification}]{KerrCTwo}. Corner removal is
\cite[\CompanionRef{signed}{test:corner-null}]{KerrCOne}; its section-volume argument
applies to any smooth positive measure on $T\Sigma$.
The geometric estimate used there is
\cite[\CompanionRef{foundation}{cmp:corner-ratios}]{KerrCTwo}. Its hypotheses are the low interior
bounds of \Cref{prop:bg-package} and a future unit timelike wedge
track starting sufficiently far in both optical coordinates.
It imposes no curvature bound or regular-extension assumption.
These statements classify the first comparison runout; that
event may precede the full geodesic endpoint in $M_d$.
\end{proof}

\subsection{Open branch tubes and launch feet}

The next lemma records uniformity in a neighborhood of a seed,
as well as the fixed-size angular footprint needed by the packet.
For a unit timelike tangent $\dot c$, write
$\pi_u=g(\dot c,\partial_u)$ and
$P_A=g(\dot c,\partial_{\theta^A})$ for its covariant momenta.
The remaining component $\pi_v$ is determined by the future unit
mass-shell equation. Thus $(u,\theta,\pi_u,P)$ are canonical state
coordinates on a transverse $v$ section.

A packet launch patch is a relatively compact coordinate patch in
$(v-T,\theta)$ on $\mathcal S$. We call it admissible when the
covectors $\dd v$ have backward null rays crossing the entry, horizon
and original initial hypersurfaces transversely, escaping into the
selected exterior end, with uniformly nondegenerate geometry on the
fixed-depth entry collar. The required fixed-order coefficient and
flow bounds are those of
\cite[\CompanionRef{foundation}{pk2:escape},
\CompanionRef{foundation}{pk2:beam-proposition}]{KerrCTwo}.
Shrinking a patch with these properties preserves them.

\begin{lemma}[A finite observation tube]\label{lem:bg-branch-tube}
Fix a branch seed, with optical ordering
$u\to u_*\in\R$, $v\to\infty$. There is a relatively compact
open neighborhood $\mathcal N$ of this seed, as small as desired, with
the following properties.
\begin{enumerate}[label=\textup{(\roman*)}]
\item All seeds in $\overline{\mathcal N}$ have an initial compact flow
tube from $\Sigma$ to a fixed late $v$ section, followed by
tracks contained in a common compact $u$ range and a sphere
coordinate chart. The section hits in the earlier flow tube are
transverse and single-valued. In canonical state variables
$(u,\theta,\pi_u,P)$, both momenta are bounded and
\begin{equation}\label{bg:p}
 p=-(\pi_u+b^AP_A)=a\dot v
\end{equation}
has positive upper and lower bounds. Each track has branch
runout in this ordering.
\item For a sufficiently small fixed interval $J$ about zero and
an unbounded sequence of $T$, there are product observation boxes
$\mathcal Q_T$ of fixed side lengths in $(u,v-T,\theta)$ containing
all hitting positions of these tracks at $v\in T+J$. The closure of
a fixed enlargement stays in the common $u$ and angular-chart margins.
The angular feet on $\mathcal S$ of the $L_0$ segments from every
point of that enlarged closure lie, with a fixed interior margin,
in one admissible small packet launch patch. The patch size is fixed by the low-order ray
and collar construction, independently of later derivative and
accuracy choices.
\item For fixed $u_2>u_*$, $c_2>0$, and sufficiently small
$c_0>0$, the neighborhood, $J$, and box sizes can be chosen so that all
entry-to-observation track portions and enlarged observation boxes fit with strict outer
margins in
\begin{equation}\label{bg:experiment-domain}
\begin{gathered}
 0\le t\le L_T=(T+u_*)/2+c_0,\\
 u\le u_2-c_0t/T,
 \qquad v\le T+c_2-c_0t/T.
\end{gathered}
\end{equation}
The slice $x$ intervals have a positive lower length, uniformly
for large $T$. No margin from the entry surface $t=0$ is
required.
\end{enumerate}
\end{lemma}

\begin{proof}
We first recall the individual branch-state argument of
\cite[\CompanionRef{signed}{bg:branch}, \CompanionRef{signed}{attach:canonical-tube}]{KerrCOne}. Let
$\Xi=\dot\theta-b\dot u$ and $P_A=\gamma_{AB}\Xi^B$.
Unit normalization gives
\begin{equation}\label{bg:unit}
 2a\dot u\dot v=1+|\Xi|_\gamma^2,
 \qquad p=-g(\dot c,L_0)>0.
\end{equation}
The low connection identities imply
$|\dot{\log p}|\le C(\dot u+|\Xi|_\gamma)$:
the shape term is controlled by
$|\Xi|_\gamma^2/p\le2\dot u$, and lapse and torsion
coefficients are bounded. On a branch, $\int\dd u<\infty$
and $\int a\,\dd v<\infty$. By \eqref{bg:unit},
\[
 \int|\Xi|_\gamma\,\dd\tau
 \le C\left(\int\dd u\int a\,\dd v\right)^{1/2}<\infty.
\]
Hence $p,p^{-1}$ are bounded, and the angular track has finite
length because $b$ is bounded and
$\dot\theta=\Xi+b\dot u$. It has a limiting point on the
sphere and eventually stays in a fixed chart.

Solving the unit mass shell for $\pi_v$ gives the reduced
$v$-time Hamiltonian
\begin{equation}\label{bg:reduced-hamiltonian}
 \mathsf H_v=-\pi_v
  =\frac{a(1+\gamma^{AB}P_AP_B)}{2p},
 \qquad
 \frac{\dd u}{\dd v}
  =\frac{a(1+\gamma^{AB}P_AP_B)}{2p^2}.
\end{equation}
Its angular Hamilton equation and the low angular coefficient
bounds give
\[
 \left|\frac{\dd P}{\dd v}\right|
 \le C(p+|P|)\frac{\dd u}{\dd v}.
\]
Gronwall in the increasing bounded coordinate $u$ bounds $P$.
It follows that $\Xi$ and $\dot u$ are bounded. In particular
$\pi_u=-p-b^AP_A$ is bounded.

Choose a compact state neighborhood of a sufficiently late
part of this track, with $p$ bounded away from zero and with
strict $u$ and angular margins. Its radius can be fixed
independently of $v$: the low first derivatives of $b$ are
uniformly bounded, so \eqref{bg:p} has a uniform state-Lipschitz
bound on a bounded momentum set. On this neighborhood, every fixed state
derivative of \eqref{bg:reduced-hamiltonian} is bounded by
$Ce^{-\kappa v}e^{o(v)}$. This follows from
\Cref{prop:bg-package}, since each derivative of $a$ retains
its factor and the state is in a compact set. The Hamiltonian
field and its Lipschitz bound are therefore integrable in $v$.
Choose the initial late section so the remaining integrated
field bound is smaller than the state margins. Gronwall then
keeps all sufficiently close states in that neighborhood for
every later $v$. Their $u$ coordinates converge, $v$ is
unbounded, and they have branch runout. Compact earlier smooth
flow, followed by one further shrinking, transports this
neighborhood to $\mathcal N\subset T\Sigma$ with the stated closure
margins. This proves \textup{(i)}.

For \textup{(ii)}, let $\Phi_v(u,\theta)$ be the angular foot
at $u=-v$ of the $L_0$ characteristic through $(u,v,\theta)$.
The angular variational equation gives
\begin{equation}\label{bg:foot-jacobian}
 \norm{D_\theta\Phi_v}+
 \norm{(D_\theta\Phi_v)^{-1}}
 \le C\exp\left(C\int_{-v}^{u}(q+g_0(s))\,\dd s\right)
 \le C.
\end{equation}
Here $\int_{-v}^{u}q\,\dd s\le q_0/\kappa$ and
\eqref{bg:cell-sum} controls the other integral. Differentiation
in the terminal $u$ contributes a bounded factor $b$.
Differentiation in $v$ contributes the moving endpoint term
$b(-v,v,\cdot)$ and an integral of $\partial_vb$ multiplied
by the bounded angular variational maps. Since
$\partial_vb=O(a)$ and $\int_{-v}^{u}a\,\dd s\le C$,
these terms are uniformly bounded as well. Thus the foot map
is uniformly Lipschitz in $(u,\theta,v)$ on the tubes in
question.

Shrink $\mathcal N$ and the width of $J$ so that the feet stay in an
arbitrarily small fixed neighborhood of the central foot at
$v=T$. The central feet lie on the compact sphere. An
unbounded subsequence therefore lies in one chart and one
small patch with interior margins. The fixed-patch ray and
collar estimates in \cite[\CompanionRef{signed}{bg:linear-packet-section}, \CompanionRef{signed}{bg:amplitude-independent}]{KerrCOne}, using
\cite[\CompanionRef{foundation}{pk2:escape}, \CompanionRef{foundation}{pk2:beam-proposition}]{KerrCTwo}, apply uniformly to all
centers in a sufficiently small low-order patch. Their patch
size is independent of later finite beam orders; higher
accuracy changes those orders and constants. The integrable
state-Lipschitz bound from \textup{(i)} makes the hitting positions
uniformly close to the central positions when $\mathcal N$ is small; the
Hamiltonian field also controls their variation across the fixed
interval $J$. Choose fixed product-box radii larger than these
variations but small relative to the launch margin. The uniform
foot-map Lipschitz bound then keeps the feet of a fixed enlargement
inside that margin. These choices are made while shrinking $\mathcal N,J$,
before any later derivative order. This supplies $\mathcal Q_T$ and proves
\textup{(ii)}.

Finally choose $c_0$ small relative to $u_2-u_*$ and $c_2$,
then shrink the state neighborhood, $J$, and box radii so that every
point of the enlarged observation boxes has strict margins from all three
upper faces in \eqref{bg:experiment-domain}. Future causal
vectors in \eqref{bg:metric} have $\dd u,\dd v\ge0$, and
$t$ is temporal. Hence each earlier point of its track from
$\mathcal S$ satisfies the same strict side inequalities:
both $u+c_0t/T$ and $v+c_0t/T$ decrease toward the past.
The slice interval has length
\[
 T+c_2+u_2-2t-2c_0t/T.
\]
At $t=L_T$ this tends to $c_2+u_2-u_*-3c_0>0$ after the
choice of $c_0$; earlier intervals are longer. This proves
\textup{(iii)}. This finite-domain fitting condition concerns
only the track up to observation; its later portion need not
remain in \eqref{bg:experiment-domain}.
\end{proof}

\subsection{Finite smooth persistence}

We also record the form of smooth Cauchy stability that will be
used. Given a fixed datum, a compact subset of its smooth
development, compact earlier flow tubes connecting it to the
initial surface, and finitely many derivative orders, nearby
smooth constraint data have corresponding neighborhoods with
metrics arbitrarily close in those norms. The identifications
come from local wave reductions, domain of dependence, and
geometric uniqueness. In particular, strict finite observation
and transversality inequalities persist. This is the compact
stability statement in \cite[\CompanionRef{foundation}{cat:closed-tests}]{KerrCTwo}, also used
in \cite[\CompanionRef{signed}{test:closed}]{KerrCOne}, based on
\cite{ChoquetBruhat1952,ChoquetBruhatGeroch1969}.

Each use is at a fixed finite experiment. The permitted data
neighborhood can depend on its time and derivative orders.
Neither persistence of a common optical gauge for nearby data
nor continuity of maximal geodesic lifetime is required.
